\section{Related Work}
\label{sec:related}

One framing prevents most confusion with prior art: we do not propose another retriever, reranker, or fusion rule, and our gains do not come from a stronger embedding model. We certify a \emph{deployment decision}---which existing physical plan to run for a workload window, and when to abstain---under a distribution-free guarantee. Read this way, the closest neighbors fall into four groups, and \CWC{} differs from each in object or granularity rather than in the retrievers it uses: database query/physical-design optimization (different operator substrate and no risk guarantee), per-query routing and selective processing (per-instance, no workload-window guarantee), risk control / conformal methods (per-instance prediction sets rather than a certified plan choice under shift), and hybrid-retrieval operators and benchmarks (the actions we select among, not selectors themselves).

\parhead{Database query optimization.}
Classical query optimizers choose physical access paths and operator orders from a cost model (System R's cost-based access-path selection \cite{selinger1979accesspath}, Volcano's extensible search \cite{graefe1993volcano}), physical-design tools recommend indexes from workload-level what-if analysis \cite{chaudhuri1998autoadmin}, and adaptive systems such as Eddies re-route during execution under runtime shifts \cite{avnur2000eddies}. \WCR{} follows this tradition but changes the operator substrate: the physical actions are hybrid retrieval plans, not relational joins or scans, and the signals are retrieval telemetry rather than only cardinality estimates.

\parhead{Workload-level and transductive query optimization.}
A recent line of database work optimizes plans at the level of a workload rather than a single query. \textsc{LimeQO} casts offline query optimization as low-rank learning over a workload, exploiting structure shared across queries \cite{limeqo2025}, and Bayesian-optimization planners learn offline plans for repeated workloads \cite{bayesqo2025}. \WCR{} shares this workload-level, transductive stance but on a different substrate (hybrid-retrieval physical plans, not relational latency) and, unlike these methods, attaches a finite-sample guarantee that remains valid under bounded workload shift. On the retrieval side, \textsc{SearchGym} proposes a compositional configuration algebra for hybrid-search orchestration \cite{searchgym2026}; that work is a benchmarking and reproducibility substrate, whereas \WCR{} is an optimizer that \emph{selects and certifies} plans over such a space under constraints.

\parhead{Risk control and selective processing for retrieval.}
Conformal risk control gives distribution-free guarantees, and recent work applies it to ranked retrieval: two-stage conformal risk control bounds retrieval and ranking risk for a single ranker's output set \cite{angelopoulos2024twostagecrc}, while selective conformal risk control first selects confident instances and then applies conformal control \cite{xu2025selectiveCRC}. Per-query selective processing chooses a search configuration per request, often using query performance prediction, but a recent study shows these per-query gains are small and unreliable across collections \cite{qpplimits2025,selectiveqp2024}. Our certificate differs in granularity and object: it controls the \emph{eval-window} violation rate of a deployed per-cell plan under covariate/workload shift, exploiting the exactly observed cell masses, rather than constructing a per-instance prediction set. This window-level framing is also what lets us avoid the per-query unreliability documented in \cite{qpplimits2025}.

\parhead{Certified plan verification (the closest prior work).}
The nearest neighbor is CP4LQO~\cite{cp4lqo2025}, which makes learned \emph{relational} query plans verifiable by using conformal prediction to bound each plan's \emph{per-query latency} before execution, with adaptive-CP handling of distribution shift and a plan search guided by those bounds. We share its goal---a deployable guarantee on a learned plan choice under shift---but differ on three axes that are important to take into account. \emph{Object:} CP4LQO certifies a per-query latency interval; we certify a \emph{workload-window SLO-violation rate} $R_Q$, the aggregate quantity an SLA or admission controller actually needs, where per-query risk does not accumulate over the window. \emph{Granularity:} CP4LQO reasons per query, so its statistical risk grows with the number of queries; we partition the window into $K$ cells and, because the eval-window cell masses $m_\cell/M$ are observed exactly (transductive), only the $K$ within-cell expectations carry statistical risk and the mass term is estimation-error-free. \emph{Domain and operation:} CP4LQO targets relational SQL and guides plan \emph{construction}; we target constrained \emph{hybrid retrieval} (BM25/dense/SPLADE/ColBERTv2/fusion/rerank) and make the certificate operational through abstention to a selection-split feasibility-best fallback. Table~\ref{tab:certtax} places both in a small design space. We position our contribution as a workload-window certified \emph{compiler} for hybrid retrieval, not a new conformal-prediction technique.

\begin{table}[t]
\centering
\caption{Design space of certified plan selection under shift. Existing certified-plan work bounds per-query latency for relational SQL; \CWC{} occupies the workload-window violation-rate cell for hybrid retrieval.}
\label{tab:certtax}
\footnotesize
\resizebox{\columnwidth}{!}{%
\begin{tabular}{lll}
\toprule
 & \textbf{Per-query latency interval} & \textbf{Window violation-rate} \\
\midrule
Relational SQL & CP4LQO~\cite{cp4lqo2025} & --- \\
Hybrid retrieval & (per-query CP, unreliable~\cite{qpplimits2025}) & \textbf{\CWC{} (ours)} \\
\bottomrule
\end{tabular}}
\end{table}

\parhead{Learned query optimization, routing, and policy selection.}
Learning-augmented optimizers such as Bao steer query optimization while integrating with existing engines \cite{marcus2021bao}, and per-query routing / offline policy selection (e.g.\ contextual-bandit action selection with offline evaluation \cite{li2010contextualbandit}) choose an action per request from request-local features or logged rewards. \WCR{} learns plan choices from prior executions like these, and we use query-only and scalar-telemetry routers as baselines, but its unit of generalization is a workload cell in retrieval-telemetry space: it observes an unlabeled window, builds reusable cells, and returns a plan frontier and risk report rather than a per-instance relational-cost prediction or contextual classification.

\parhead{Hybrid-retrieval operators, reranking, and benchmarks.}
The actions we select among are standard hybrid-retrieval operators: BM25 \cite{robertson2009bm25}, dense (E5 \cite{wang2022e5}, BGE-M3 \cite{chen2024m3embedding}, Qwen3-Embedding \cite{zhang2025qwen3embedding}, NV-Embed \cite{lee2024nvembed}), learned-sparse SPLADE \cite{formal2021splade}, late-interaction ColBERT/ColBERTv2 \cite{khattab2020colbert,santhanam2021colbertv2}, reciprocal-rank fusion \cite{cormack2009rrf}, and neural rerankers such as RankZephyr \cite{pradeep2023rankzephyr}, which trade quality for latency. \WCR{} treats each as a physical action that may or may not be chosen for a cell, rather than proposing a new operator. Finally, the BEIR benchmark systematized the zero-shot evaluation of retrieval models \cite{thakur2021beir}; its metrics, however, score ranking quality alone and say nothing about deployability under a latency SLO; we evaluate constrained score, satisfaction, latency, regret, and rank-gated comparisons so quality and systems feasibility are tested together.
